% Осень26!A5; нумерация внутри столбца без пустых строк.
% Источник: лекция 1, §6 (строка 146)
% Источник: лекция 1, §8 (строка 209)
\ticket{03}{Лемма о дифференцировании ``сдвига\textquotesingle\textquotesingle{}}
\label{ticket:03}

\subsection*{Мультииндекс}

\textbf{Мотивировка.} Дважды лектор «выкручивался», ограничиваясь двумерной ситуацией, «но сколько же можно»: для формулы Тейлора в произвольной размерности заводится новый персонаж --- \textbf{мультииндекс} («мультяшный индекс»).

\begin{definition}

Мультииндекс --- это набор \[k=(k_1,k_2,\dots,k_m)\in\bigl(\mathbb{N}\cup\{0\}\bigr)^m,\] т.е. «векторнозначный целочисленный индекс с \(m\) координатами», каждая координата \(k_i\) --- неотрицательное целое число.

Связанные обозначения (для \(x\in\mathbb{R}^m\), \(f\) --- числовая функция):

\end{definition}

\textbf{Высота мультииндекса:} \(|k|=k_1+\cdots+k_m\).

\textbf{Факториал мультииндекса:} \(k!=k_1!\cdots k_m!\).

\textbf{Степень вектора:} \(x^k=x_1^{k_1}\cdots x_m^{k_m}\).

\textbf{Производная по мультииндексу:} \[D^k f=f^{(k)}=\frac{\partial^{|k|}f}{\partial x_1^{k_1}\cdots\partial x_m^{k_m}}.\]

Пояснения.

\begin{itemize}
\tightlist
\item
  Слово «норма» для \(\lvert k\rvert\) не используется --- оно «занято» во многих местах; принят термин \textbf{высота}.
\item
  \(x^k\) определено при любых \(x\): степени целые неотрицательные, «нет проблем типа корня из \(-1\)»; это одночлен --- произведение степеней координат.
\item
  \(D^k f\): по переменной \(x_1\) дифференцируем \(k_1\) раз, по \(x_2\) --- \(k_2\) раз, и т.д.; всего \(\lvert k\rvert\) дифференцирований. Мультииндекс описывает «комплект переменных с кратностями»: \(x_1\) с кратностью \(k_1\), \ldots, \(x_m\) с кратностью \(k_m\) (кратности могут быть нулевыми --- тогда некоторых переменных нет). Корректность (независимость от порядка) гарантирует теорема §5; определением пользоваться можно при \(f\in C^r(E)\), \(r\ge\lvert k\rvert\).
\item
  \textbf{Откуда обозначение} \(\partial^2/\partial x_1^2\): «когда вы дифференцируете по \(x_1\), вы применяете оператор \(\frac{\partial}{\partial x_1}\); если ещё раз --- применяете его снова; любой математик, увидев два таких значка рядом, скажет: \(d\cdot d=d^2\), \(\partial x_1\cdot\partial x_1=\partial x_1^2\)» (скобки обычно не ставят). Обозначение и возникло из этой «наивной» алгебры операторов.
\end{itemize}

\subsection*{Лемма о дифференцировании «сдвига» (лемма 2)}

У нас появилась «новейшая технология дифференцирования высокого порядка» --- упражнение на её применение: как дифференцируются сдвиги (сужения функции на прямые).

\begin{lemma}

\emph{2 (о дифференцировании «сдвига».} Пусть \(E\subset\mathbb{R}^m\) открыто, \(f\in C^r(E)\), точка \(a\in E\), вектор \(h\in\mathbb{R}^m\), и пусть \(a+th\in E\) для всех \(t\in(-1,1)\) (т.е. вдоль направления \(h\) точка \(a\) «с запасом внутри» \(E\); ноль глубоко внутри допустимого интервала). Определим функцию одной переменной \[\varphi(t):=f(a+th),\qquad t\in(-1,1).\] Тогда для каждого \(k\le r\) и всех \(t\in(-1,1)\) \[\varphi^{(k)}(t)=\sum_{|g|=k}\frac{k!}{g!}\;h^{g}\;D^{g}f\,(a+th),\] где суммирование --- по всем мультииндексам \(g\) высоты \(k\).

\textbf{Требует сверки.} Оговорка об области \(t\): в лекции требовалось \(a+th\in E\) «для любого \(t\) из отрезка \([-1,1]\)», с ремаркой, что это «окрестность нуля и важно лишь, что ноль глубоко внутри»; в справочнике --- интервал \((-1,1)\). Расхождение несущественно: нужна лишь окрестность нуля с запасом, содержащая отрезок \([0,1]\) после замены \(x=a+h\).

\end{lemma}

\textbf{Идея доказательства.} Дифференцирование \(\varphi\) по \(t\) --- это производная композиции: по правилу «дифференцировать по каждому вхождению буквы \(t\)» первое же дифференцирование даёт сумму \(\sum_i f'_{x_i}(a+th)\,h_i\); каждое следующее дифференцирование «выбрасывает» ещё одну сумму, очередной множитель \(h_i\) и очередную «стрелочку» производной у \(f\). После \(k\) дифференцирований получается \(k\) вложенных сумм по «маршрутам» \((i_1,\dots,i_k)\) выбора переменных. Одночлен-мультипликатор \(h_{i_1}\cdots h_{i_k}=h^g\) зависит только от того, \textbf{сколько раз} каждая переменная выбрана (т.е. от мультииндекса \(g\)), и по \textbf{полиномиальной формуле} (лемма 1) каждый \(h^g\) встречается ровно \(k!/g!\) раз; а по \textbf{теореме о независимости от порядка} (§2, §5) все производные с одним и тем же мультииндексом \(g\) равны. Приведение подобных членов и даёт формулу.

\begin{proof}

\textbf{Доказательство} (после перерыва). Выпишем покоординатно \(\varphi(t)=f(a_1+th_1,\,a_2+th_2,\dots,a_m+th_m)\) --- видно, где прячется переменная \(t\).

\emph{Первое дифференцирование.} По правилу дифференцирования сложной функции (правило «по каждому вхождению \(t\) независимо», доказанное во 2 семестре после теоремы о производной композиции): вхождение \(t\) на \(i\)-м месте --- это дифференцирование \(f\) по \(i\)-й переменной, умноженное на производную \(a_i+th_i\) по \(t\), т.е. на \(h_i\): \[\varphi'(t)=\sum_{i=1}^{m}f'_{x_i}(a+th)\,h_i.\]

\emph{Второе дифференцирование.} Дифференцируем каждое слагаемое снова по \(t\); появляется вторая сумма: \[\varphi''(t)=\sum_{i=1}^{m}\sum_{j=1}^{m}f''_{x_jx_i}(a+th)\,h_i h_j\] (новое дифференцирование дописывается слева --- по нашему соглашению; «теперь я могу не думать о порядке --- теорема говорит, что без разницы»).

\emph{Общий шаг (индукция).} \(k\) дифференцирований дают \(k\) вложенных сумм: \[\varphi^{(k)}(t)=\sum_{i_1=1}^{m}\sum_{i_2=1}^{m}\dots\sum_{i_k=1}^{m}\frac{\partial^{k}f}{\partial x_{i_1}\partial x_{i_2}\cdots\partial x_{i_k}}(a+th)\;\cdot\;h_{i_1}h_{i_2}\cdots h_{i_k}.\] Алгоритм прозрачен: «каждое дифференцирование выбрасывает сумму, очередную производную и очередной множитель».

\emph{Приведение подобных.} Зададим мультииндекс \(g\) счётчиками: \(g_s\) --- сколько раз среди \(i_1,\dots,i_k\) встречается номер \(s\). Тогда \(h_{i_1}\cdots h_{i_k}=h^g\), \(|g|=k\), и фиксированный одночлен \(h^g\) возникает при всех «перестановках» его индексов. По лемме 1 (полиномиальной формуле) таких наборов-перестановок ровно \(k!/g!\). Все соответствующие производные равны между собой по теореме о независимости от порядка дифференцирования (они различаются лишь порядком одних и тех же дифференцирований). Значит, суммируя, получаем \[\varphi^{(k)}(t)=\sum_{|g|=k}\frac{k!}{g!}\,h^{g}\,D^{g}f\,(a+th).\]

\end{proof}
